% comparison_FR.tex — WAC confronted with KAPE and Eric Zimmerman's tools, French edition.
% The English edition is comparison_EN.tex, in this folder: both are maintained together,
% section for section.
%
% Build (XeLaTeX, two passes for the cover page and the table of contents):
%     latexmk -xelatex comparison_FR.tex
% The cover page is ../cover_comparison_FR.png, derived from ../cover_FR.png by
% ../retitle-cover.py. The layout is ../wac-style.tex, shared with every WAC document.
%
% Every figure comes from backlog/B-015 and the runs of vmtest/ez-compare.sh and
% vmtest/kape-compare.sh, which reproduce them.

\documentclass[11pt,a4paper,oneside]{report}

% ---------------------------------------------------------------- language and labels
\usepackage[french]{babel}
\newcommand{\wacdoctitle}{Face à KAPE et aux outils EZ}
\newcommand{\wacpageof}{sur}
\newcommand{\wacnotetitle}{À retenir}
\newcommand{\wacwarningtitle}{Limite}
% Layout shared with every WAC document.
% wac-style.tex — layout shared by every WAC document (user manuals, testing
% guides), French and English editions alike: one copy, so that they cannot
% drift apart.
%
% A document sets, BEFORE % wac-style.tex — layout shared by every WAC document (user manuals, testing
% guides), French and English editions alike: one copy, so that they cannot
% drift apart.
%
% A document sets, BEFORE % wac-style.tex — layout shared by every WAC document (user manuals, testing
% guides), French and English editions alike: one copy, so that they cannot
% drift apart.
%
% A document sets, BEFORE \input{../wac-style.tex}:
%   \usepackage[<language>]{babel}
%   \newcommand{\wacdoctitle}{...}      running header: "WAC — <title>"
%   \newcommand{\wacpageof}{...}        footer: "Page N <of> M"
%   \newcommand{\wacnotetitle}{...}     default title of the "note" box
%   \newcommand{\wacwarningtitle}{...}  default title of the "attention" box
% and AFTER it its \hypersetup (PDF title, author, subject).

% ---------------------------------------------------------------- fonts, language
\usepackage{fontspec}
\setmainfont{Noto Sans}
\setmonofont{DejaVu Sans Mono}[Scale=0.84]

% ---------------------------------------------------------------- layout
\usepackage[top=26mm,bottom=26mm,left=24mm,right=22mm,headheight=14pt]{geometry}
\usepackage{graphicx}
\graphicspath{{../}}
\usepackage{xcolor}
\usepackage{tikz}
\usepackage{booktabs,tabularx,longtable,array}
\usepackage{enumitem}
\usepackage{fancyhdr}
\usepackage{titlesec}
\usepackage[most]{tcolorbox}
\usepackage{listings}
\usepackage{url}
\usepackage{lastpage}
\usepackage[hidelinks,unicode,pdfpagelabels]{hyperref}

% Colours sampled from the cover page
\definecolor{marine}{HTML}{0A1A33}
\definecolor{bleu}{HTML}{2F6DB5}
\definecolor{accent}{HTML}{1F66D1}
\definecolor{fond}{HTML}{EEF3FA}
\definecolor{alerte}{HTML}{B3261E}
\definecolor{fondalerte}{HTML}{FBEDEC}
\definecolor{gris}{HTML}{5B6573}

\setlength{\parindent}{0pt}
\setlength{\parskip}{0.55em}
\linespread{1.08}
\setlist{itemsep=0.25em,topsep=0.35em}
\urlstyle{tt}
\renewcommand{\arraystretch}{1.25}
% Left-aligned table columns: justified narrow columns opened wide gaps
% between words.
\newcolumntype{L}[1]{>{\raggedright\arraybackslash}p{#1}}
% Monospaced file and field names cannot be hyphenated: let the line's spaces
% stretch instead.
\setlength{\emergencystretch}{3em}

% Headings
\titleformat{\chapter}[display]
  {\color{marine}\bfseries\raggedright}
  {\color{bleu}\large\MakeUppercase{\chaptertitlename}~\thechapter}
  {0.3em}
  {\Huge}[{\color{accent}\titlerule[2pt]}]
\titlespacing*{\chapter}{0pt}{0pt}{1.6em}
\titleformat{\section}{\color{marine}\Large\bfseries}{\color{bleu}\thesection}{0.8em}{}
\titleformat{\subsection}{\color{marine}\large\bfseries}{\color{bleu}\thesubsection}{0.7em}{}
\titleformat{\subsubsection}{\color{bleu}\normalsize\bfseries}{}{0em}{}

% Headers and footers
\pagestyle{fancy}
\fancyhf{}
\fancyhead[L]{\small\color{gris}WAC — \wacdoctitle}
\fancyhead[R]{\small\color{gris}\leftmark}
% Footer: "Page N of M" (the word set by \wacpageof), readable. A bare, small, light-grey number went
% unnoticed.
\newcommand{\piedpage}{\small\color{marine}\textbf{Page \thepage}\ \wacpageof\ \pageref*{LastPage}}
\fancyfoot[C]{\piedpage}
\renewcommand{\footrulewidth}{0.4pt}
\renewcommand{\footrule}{\hbox to\headwidth{\color{fond}\leaders\hrule height \footrulewidth\hfill}}
\renewcommand{\headrulewidth}{0.4pt}
\renewcommand{\headrule}{\hbox to\headwidth{\color{fond}\leaders\hrule height \headrulewidth\hfill}}
\renewcommand{\chaptermark}[1]{\markboth{#1}{}}
\fancypagestyle{plain}{\fancyhf{}\fancyfoot[C]{\piedpage}\renewcommand{\headrulewidth}{0pt}}

% Boxes
\newtcolorbox{note}[1][\wacnotetitle]{enhanced,breakable,colback=fond,colframe=bleu,
  boxrule=0pt,leftrule=3pt,arc=0pt,left=8pt,right=8pt,top=6pt,bottom=6pt,
  title={\bfseries #1},coltitle=bleu,colbacktitle=fond,attach title to upper={\par\smallskip}}
\newtcolorbox{attention}[1][\wacwarningtitle]{enhanced,breakable,colback=fondalerte,colframe=alerte,
  boxrule=0pt,leftrule=3pt,arc=0pt,left=8pt,right=8pt,top=6pt,bottom=6pt,
  title={\bfseries #1},coltitle=alerte,colbacktitle=fondalerte,attach title to upper={\par\smallskip}}

% Code
\lstset{basicstyle=\ttfamily\small,columns=fullflexible,keepspaces=true,
  backgroundcolor=\color{fond},frame=l,framerule=3pt,rulecolor=\color{bleu},
  xleftmargin=9pt,framexleftmargin=6pt,aboveskip=0.8em,belowskip=0.8em,
  breaklines=true,breakatwhitespace=false,showstringspaces=false,
  literate={é}{{\'e}}1 {è}{{\`e}}1 {à}{{\`a}}1 {ê}{{\^e}}1 {ô}{{\^o}}1 {ç}{{\c c}}1
           {É}{{\'E}}1 {—}{{---}}1 {’}{{'}}1 {…}{{...}}1 {«}{{\guillemotleft}}1 {»}{{\guillemotright}}1}

% Shortcuts
\newcommand{\wac}{\textsc{wac}}
\newcommand{\opt}[1]{\texttt{-{}-#1}}
\newcommand{\champ}[1]{\texttt{#1}}
\newcommand{\fichier}[1]{\texttt{#1}}

% ============================================================ COVER PAGE
% The cover images are 2:3 (1024 x 1536), taller than an A4 page once widened
% to it: centred, they lost 9 mm at the top — the CEC logo — and 9 mm at the
% bottom. Raised 4 mm above the top of the page instead: the logo, which starts
% 6 mm below the image's edge, stays whole, and the bottom row of text keeps
% 7 mm to the page's edge; only a decorative corner below it is cut.
% Usage: \coverpage{cover_FR.png}
\newcommand{\coverpage}[1]{%
  \begin{titlepage}
  \thispagestyle{empty}
  \begin{tikzpicture}[remember picture,overlay]
    \node[inner sep=0pt,anchor=north,yshift=4mm] at (current page.north)
      {\includegraphics[width=\paperwidth]{#1}};
  \end{tikzpicture}
  \end{titlepage}}
:
%   \usepackage[<language>]{babel}
%   \newcommand{\wacdoctitle}{...}      running header: "WAC — <title>"
%   \newcommand{\wacpageof}{...}        footer: "Page N <of> M"
%   \newcommand{\wacnotetitle}{...}     default title of the "note" box
%   \newcommand{\wacwarningtitle}{...}  default title of the "attention" box
% and AFTER it its \hypersetup (PDF title, author, subject).

% ---------------------------------------------------------------- fonts, language
\usepackage{fontspec}
\setmainfont{Noto Sans}
\setmonofont{DejaVu Sans Mono}[Scale=0.84]

% ---------------------------------------------------------------- layout
\usepackage[top=26mm,bottom=26mm,left=24mm,right=22mm,headheight=14pt]{geometry}
\usepackage{graphicx}
\graphicspath{{../}}
\usepackage{xcolor}
\usepackage{tikz}
\usepackage{booktabs,tabularx,longtable,array}
\usepackage{enumitem}
\usepackage{fancyhdr}
\usepackage{titlesec}
\usepackage[most]{tcolorbox}
\usepackage{listings}
\usepackage{url}
\usepackage{lastpage}
\usepackage[hidelinks,unicode,pdfpagelabels]{hyperref}

% Colours sampled from the cover page
\definecolor{marine}{HTML}{0A1A33}
\definecolor{bleu}{HTML}{2F6DB5}
\definecolor{accent}{HTML}{1F66D1}
\definecolor{fond}{HTML}{EEF3FA}
\definecolor{alerte}{HTML}{B3261E}
\definecolor{fondalerte}{HTML}{FBEDEC}
\definecolor{gris}{HTML}{5B6573}

\setlength{\parindent}{0pt}
\setlength{\parskip}{0.55em}
\linespread{1.08}
\setlist{itemsep=0.25em,topsep=0.35em}
\urlstyle{tt}
\renewcommand{\arraystretch}{1.25}
% Left-aligned table columns: justified narrow columns opened wide gaps
% between words.
\newcolumntype{L}[1]{>{\raggedright\arraybackslash}p{#1}}
% Monospaced file and field names cannot be hyphenated: let the line's spaces
% stretch instead.
\setlength{\emergencystretch}{3em}

% Headings
\titleformat{\chapter}[display]
  {\color{marine}\bfseries\raggedright}
  {\color{bleu}\large\MakeUppercase{\chaptertitlename}~\thechapter}
  {0.3em}
  {\Huge}[{\color{accent}\titlerule[2pt]}]
\titlespacing*{\chapter}{0pt}{0pt}{1.6em}
\titleformat{\section}{\color{marine}\Large\bfseries}{\color{bleu}\thesection}{0.8em}{}
\titleformat{\subsection}{\color{marine}\large\bfseries}{\color{bleu}\thesubsection}{0.7em}{}
\titleformat{\subsubsection}{\color{bleu}\normalsize\bfseries}{}{0em}{}

% Headers and footers
\pagestyle{fancy}
\fancyhf{}
\fancyhead[L]{\small\color{gris}WAC — \wacdoctitle}
\fancyhead[R]{\small\color{gris}\leftmark}
% Footer: "Page N of M" (the word set by \wacpageof), readable. A bare, small, light-grey number went
% unnoticed.
\newcommand{\piedpage}{\small\color{marine}\textbf{Page \thepage}\ \wacpageof\ \pageref*{LastPage}}
\fancyfoot[C]{\piedpage}
\renewcommand{\footrulewidth}{0.4pt}
\renewcommand{\footrule}{\hbox to\headwidth{\color{fond}\leaders\hrule height \footrulewidth\hfill}}
\renewcommand{\headrulewidth}{0.4pt}
\renewcommand{\headrule}{\hbox to\headwidth{\color{fond}\leaders\hrule height \headrulewidth\hfill}}
\renewcommand{\chaptermark}[1]{\markboth{#1}{}}
\fancypagestyle{plain}{\fancyhf{}\fancyfoot[C]{\piedpage}\renewcommand{\headrulewidth}{0pt}}

% Boxes
\newtcolorbox{note}[1][\wacnotetitle]{enhanced,breakable,colback=fond,colframe=bleu,
  boxrule=0pt,leftrule=3pt,arc=0pt,left=8pt,right=8pt,top=6pt,bottom=6pt,
  title={\bfseries #1},coltitle=bleu,colbacktitle=fond,attach title to upper={\par\smallskip}}
\newtcolorbox{attention}[1][\wacwarningtitle]{enhanced,breakable,colback=fondalerte,colframe=alerte,
  boxrule=0pt,leftrule=3pt,arc=0pt,left=8pt,right=8pt,top=6pt,bottom=6pt,
  title={\bfseries #1},coltitle=alerte,colbacktitle=fondalerte,attach title to upper={\par\smallskip}}

% Code
\lstset{basicstyle=\ttfamily\small,columns=fullflexible,keepspaces=true,
  backgroundcolor=\color{fond},frame=l,framerule=3pt,rulecolor=\color{bleu},
  xleftmargin=9pt,framexleftmargin=6pt,aboveskip=0.8em,belowskip=0.8em,
  breaklines=true,breakatwhitespace=false,showstringspaces=false,
  literate={é}{{\'e}}1 {è}{{\`e}}1 {à}{{\`a}}1 {ê}{{\^e}}1 {ô}{{\^o}}1 {ç}{{\c c}}1
           {É}{{\'E}}1 {—}{{---}}1 {’}{{'}}1 {…}{{...}}1 {«}{{\guillemotleft}}1 {»}{{\guillemotright}}1}

% Shortcuts
\newcommand{\wac}{\textsc{wac}}
\newcommand{\opt}[1]{\texttt{-{}-#1}}
\newcommand{\champ}[1]{\texttt{#1}}
\newcommand{\fichier}[1]{\texttt{#1}}

% ============================================================ COVER PAGE
% The cover images are 2:3 (1024 x 1536), taller than an A4 page once widened
% to it: centred, they lost 9 mm at the top — the CEC logo — and 9 mm at the
% bottom. Raised 4 mm above the top of the page instead: the logo, which starts
% 6 mm below the image's edge, stays whole, and the bottom row of text keeps
% 7 mm to the page's edge; only a decorative corner below it is cut.
% Usage: \coverpage{cover_FR.png}
\newcommand{\coverpage}[1]{%
  \begin{titlepage}
  \thispagestyle{empty}
  \begin{tikzpicture}[remember picture,overlay]
    \node[inner sep=0pt,anchor=north,yshift=4mm] at (current page.north)
      {\includegraphics[width=\paperwidth]{#1}};
  \end{tikzpicture}
  \end{titlepage}}
:
%   \usepackage[<language>]{babel}
%   \newcommand{\wacdoctitle}{...}      running header: "WAC — <title>"
%   \newcommand{\wacpageof}{...}        footer: "Page N <of> M"
%   \newcommand{\wacnotetitle}{...}     default title of the "note" box
%   \newcommand{\wacwarningtitle}{...}  default title of the "attention" box
% and AFTER it its \hypersetup (PDF title, author, subject).

% ---------------------------------------------------------------- fonts, language
\usepackage{fontspec}
\setmainfont{Noto Sans}
\setmonofont{DejaVu Sans Mono}[Scale=0.84]

% ---------------------------------------------------------------- layout
\usepackage[top=26mm,bottom=26mm,left=24mm,right=22mm,headheight=14pt]{geometry}
\usepackage{graphicx}
\graphicspath{{../}}
\usepackage{xcolor}
\usepackage{tikz}
\usepackage{booktabs,tabularx,longtable,array}
\usepackage{enumitem}
\usepackage{fancyhdr}
\usepackage{titlesec}
\usepackage[most]{tcolorbox}
\usepackage{listings}
\usepackage{url}
\usepackage{lastpage}
\usepackage[hidelinks,unicode,pdfpagelabels]{hyperref}

% Colours sampled from the cover page
\definecolor{marine}{HTML}{0A1A33}
\definecolor{bleu}{HTML}{2F6DB5}
\definecolor{accent}{HTML}{1F66D1}
\definecolor{fond}{HTML}{EEF3FA}
\definecolor{alerte}{HTML}{B3261E}
\definecolor{fondalerte}{HTML}{FBEDEC}
\definecolor{gris}{HTML}{5B6573}

\setlength{\parindent}{0pt}
\setlength{\parskip}{0.55em}
\linespread{1.08}
\setlist{itemsep=0.25em,topsep=0.35em}
\urlstyle{tt}
\renewcommand{\arraystretch}{1.25}
% Left-aligned table columns: justified narrow columns opened wide gaps
% between words.
\newcolumntype{L}[1]{>{\raggedright\arraybackslash}p{#1}}
% Monospaced file and field names cannot be hyphenated: let the line's spaces
% stretch instead.
\setlength{\emergencystretch}{3em}

% Headings
\titleformat{\chapter}[display]
  {\color{marine}\bfseries\raggedright}
  {\color{bleu}\large\MakeUppercase{\chaptertitlename}~\thechapter}
  {0.3em}
  {\Huge}[{\color{accent}\titlerule[2pt]}]
\titlespacing*{\chapter}{0pt}{0pt}{1.6em}
\titleformat{\section}{\color{marine}\Large\bfseries}{\color{bleu}\thesection}{0.8em}{}
\titleformat{\subsection}{\color{marine}\large\bfseries}{\color{bleu}\thesubsection}{0.7em}{}
\titleformat{\subsubsection}{\color{bleu}\normalsize\bfseries}{}{0em}{}

% Headers and footers
\pagestyle{fancy}
\fancyhf{}
\fancyhead[L]{\small\color{gris}WAC — \wacdoctitle}
\fancyhead[R]{\small\color{gris}\leftmark}
% Footer: "Page N of M" (the word set by \wacpageof), readable. A bare, small, light-grey number went
% unnoticed.
\newcommand{\piedpage}{\small\color{marine}\textbf{Page \thepage}\ \wacpageof\ \pageref*{LastPage}}
\fancyfoot[C]{\piedpage}
\renewcommand{\footrulewidth}{0.4pt}
\renewcommand{\footrule}{\hbox to\headwidth{\color{fond}\leaders\hrule height \footrulewidth\hfill}}
\renewcommand{\headrulewidth}{0.4pt}
\renewcommand{\headrule}{\hbox to\headwidth{\color{fond}\leaders\hrule height \headrulewidth\hfill}}
\renewcommand{\chaptermark}[1]{\markboth{#1}{}}
\fancypagestyle{plain}{\fancyhf{}\fancyfoot[C]{\piedpage}\renewcommand{\headrulewidth}{0pt}}

% Boxes
\newtcolorbox{note}[1][\wacnotetitle]{enhanced,breakable,colback=fond,colframe=bleu,
  boxrule=0pt,leftrule=3pt,arc=0pt,left=8pt,right=8pt,top=6pt,bottom=6pt,
  title={\bfseries #1},coltitle=bleu,colbacktitle=fond,attach title to upper={\par\smallskip}}
\newtcolorbox{attention}[1][\wacwarningtitle]{enhanced,breakable,colback=fondalerte,colframe=alerte,
  boxrule=0pt,leftrule=3pt,arc=0pt,left=8pt,right=8pt,top=6pt,bottom=6pt,
  title={\bfseries #1},coltitle=alerte,colbacktitle=fondalerte,attach title to upper={\par\smallskip}}

% Code
\lstset{basicstyle=\ttfamily\small,columns=fullflexible,keepspaces=true,
  backgroundcolor=\color{fond},frame=l,framerule=3pt,rulecolor=\color{bleu},
  xleftmargin=9pt,framexleftmargin=6pt,aboveskip=0.8em,belowskip=0.8em,
  breaklines=true,breakatwhitespace=false,showstringspaces=false,
  literate={é}{{\'e}}1 {è}{{\`e}}1 {à}{{\`a}}1 {ê}{{\^e}}1 {ô}{{\^o}}1 {ç}{{\c c}}1
           {É}{{\'E}}1 {—}{{---}}1 {’}{{'}}1 {…}{{...}}1 {«}{{\guillemotleft}}1 {»}{{\guillemotright}}1}

% Shortcuts
\newcommand{\wac}{\textsc{wac}}
\newcommand{\opt}[1]{\texttt{-{}-#1}}
\newcommand{\champ}[1]{\texttt{#1}}
\newcommand{\fichier}[1]{\texttt{#1}}

% ============================================================ COVER PAGE
% The cover images are 2:3 (1024 x 1536), taller than an A4 page once widened
% to it: centred, they lost 9 mm at the top — the CEC logo — and 9 mm at the
% bottom. Raised 4 mm above the top of the page instead: the logo, which starts
% 6 mm below the image's edge, stays whole, and the bottom row of text keeps
% 7 mm to the page's edge; only a decorative corner below it is cut.
% Usage: \coverpage{cover_FR.png}
\newcommand{\coverpage}[1]{%
  \begin{titlepage}
  \thispagestyle{empty}
  \begin{tikzpicture}[remember picture,overlay]
    \node[inner sep=0pt,anchor=north,yshift=4mm] at (current page.north)
      {\includegraphics[width=\paperwidth]{#1}};
  \end{tikzpicture}
  \end{titlepage}}


\hypersetup{pdftitle={WAC — Face à KAPE et aux outils EZ},
  pdfauthor={Centre d'Excellence Cyberdéfense Aérospatiale (CEC)},
  pdfsubject={WAC confronté à KAPE et aux outils d'Eric Zimmerman : temps, précision, formats}}

\begin{document}

% ============================================================ COVER PAGE
\pagenumbering{roman}
\coverpage{cover_comparison_FR.png}

% ============================================================ INFORMATION PAGE
\thispagestyle{empty}
\vspace*{\fill}
{\color{marine}\Large\bfseries WAC — Windows Artefacts Collector}\par
{\color{bleu}Face à KAPE et aux outils d'Eric Zimmerman}\par
\medskip
\begin{tabular}{@{}ll@{}}
Version du logiciel & 1.4.0 \\
Édition du document & octobre 2026 \\
Éditeur & Centre d'Excellence Cyberdéfense Aérospatiale (CEC) \\
        & École de l'air et de l'espace, base aérienne 701, Salon-de-Provence \\
\end{tabular}

\bigskip
{\small\color{gris}Ce document confronte \wac{} aux outils de référence de la
communauté : KAPE pour la collecte, les outils en ligne de commande d'Eric
Zimmerman pour le décodage. Sur les mêmes pièces d'une même machine, il compare
les temps d'exécution, la justesse des valeurs --- chaque écart étant tranché par
une référence indépendante --- et les formats de sortie. Il dit ce que la
confrontation a appris à \wac{}, et ce qu'elle lui reproche encore.}
\vspace*{2cm}
\clearpage

\pagenumbering{arabic}
\tableofcontents

% ============================================================ 1
\chapter{La démarche}

\section{Ce qui est comparé}

Deux questions, deux comparaisons.

\begin{description}[style=nextline,leftmargin=1em]
  \item[Collecte : \wac{} face à KAPE] Les deux outils collectent la même
        machine l'un après l'autre ; on compare ce que chacun copie, le temps
        qu'il y met, et l'empreinte SHA-1 des fichiers qu'ils copient tous deux.
  \item[Décodage : \wac{} face aux outils EZ] Sur les \emph{mêmes pièces} --- le
        magasin de pièces d'une collecte \wac{} ---, \wac{} convertit un artefact
        et l'outil d'Eric Zimmerman correspondant le décode ; les deux sorties
        sont comparées enregistrement par enregistrement et champ par champ.
\end{description}

\section{Qui a raison quand ils diffèrent}

Un écart entre deux outils ne dit pas lequel se trompe. Chaque écart a donc été
tranché par une troisième source, indépendante des deux :

\begin{itemize}[nosep]
  \item Windows lui-même : le XML qu'il rend pour un événement
        (\texttt{Get-WinEvent}), la date d'un fichier, l'identifiant d'objet
        NTFS (\texttt{fsutil objectid}), le registre en ligne, l'empreinte d'un
        fichier (\texttt{Get-FileHash}) ;
  \item une autre implémentation : \texttt{dissect.esedb} (Fox-IT) pour les bases
        ESE ;
  \item le format lui-même : numéros de séquence et sommes Marvin32 des journaux
        de ruche, lus octet par octet ;
  \item le recoupement de deux artefacts de la même machine : le Prefetch pour
        l'exécution, l'horodatage d'une ligne pour ses autres dates.
\end{itemize}

Chaque écart est ensuite classé : \emph{défaut de \wac{}} (corrigé), \emph{erreur
de l'outil EZ} (vérifiée), ou \emph{convention} de l'un des deux --- deux formes
d'une même valeur.

\section{Le banc d'essai}

\begin{tabular}{@{}ll@{}}
Machine & Windows 11 Professionnel (build 26200), machine virtuelle \\
Outils EZ & version 2026.5.0 (.NET 9.0.20) \\
KAPE & version 1.3.0.2, cible \texttt{KapeTriage} \\
\wac{} & version 1.4.0 et ses corrections issues de cette confrontation \\
\end{tabular}

Tout se rejoue : \fichier{vmtest/ez-compare.sh <artefact>} chronomètre \wac{}
et l'outil EZ sur les mêmes pièces puis compare leurs sorties ;
\fichier{vmtest/kape-compare.sh} fait de même pour la collecte. Les outils EZ ont
été lancés avec leur pleine précision de date (\opt{mp}, ou \opt{dt} au format
à sept décimales quand \opt{mp} ne suffit pas), pour ne pas leur imputer une
perte de précision qu'un réglage évite.

% ============================================================ 2
\chapter{Synthèse}

\begin{longtable}{@{}L{3.2cm}L{2.4cm}L{4.5cm}L{5.0cm}@{}}
\toprule
\textbf{Artefact} & \textbf{Volume} & \textbf{Résultat} & \textbf{Erreurs de l'outil EZ, vérifiées} \\
\midrule\endhead
\$MFT --- MFTECmd & 362\,081 enreg. & chemins, états, 8 dates identiques & tailles de 1\,490 fichiers WOF et 3 journaux ouverts : Windows donne celles de \wac{} \\
EVTX --- EvtxECmd & 190\,613 évén. & champs système identiques ; valeurs expliquées & 6 valeurs fausses (Windows) ; 335 événements Perflib illisibles ; journaux vidés ignorés \\
Prefetch --- PECmd & 379 fichiers & tout identique & empreinte privée de son zéro de tête (19) \\
Raccourcis --- LECmd & 3 fichiers & identiques au 100\,ns & --- \\
Jump lists --- JLECmd & 13 raccourcis & identiques au 100\,ns & --- \\
Shellbags --- SBECmd & 27 dossiers & identiques & NodeSlot faux au premier niveau (7) ; HasExplored toujours faux \\
Amcache --- AmcacheParser & 1\,258 entrées & 124 colonnes identiques & 2\,340 valeurs inventées ; SHA-1 partiel non signalé \\
ShimCache --- AppCompat\-Cache\-Parser & 232 entrées & identiques & indicateur d'exécution trompeur \\
Corbeille --- RBCmd & 19 fichiers & identiques & --- \\
SRUM --- SrumECmd & 60\,804 lignes & identiques & 17 lignes d'énergie fausses (dissect) \\
Registre --- RECmd & 9 artefacts & identiques & 4 ruches ramenées à un état ancien \\
Collecte --- KAPE & 1\,173 fichiers communs & SHA-1 identiques & --- \\
\bottomrule
\end{longtable}

\begin{note}
Sur tout ce que les deux côtés lisent, \wac{} donne les mêmes valeurs que les
outils de référence, ou des valeurs plus justes, établies à chaque fois par une
source indépendante. La confrontation a aussi révélé 23 défauts de \wac{}, tous
corrigés (chapitre~\ref{chap:defauts}), et des lacunes de couverture, ouvertes
au suivi (section~\ref{sec:couverture}).
\end{note}

% ============================================================ 3
\chapter{Temps d'exécution}

\section{Décodage}

Même machine, mêmes pièces ; chaque outil ne traite que son artefact.

\begin{longtable}{@{}L{3.6cm}L{2.2cm}L{2.2cm}L{7.0cm}@{}}
\toprule
\textbf{Artefact} & \textbf{\wac{}} & \textbf{Outil EZ} & \textbf{Ce que fait \wac{} en plus} \\
\midrule\endhead
\$MFT (380\,Mo) & 7,3\,s & 8,1\,s & deux volumes (C: et G:) contre un seul \\
EVTX (305 journaux) & 16,8\,s & 22,0\,s & message en clair reconstitué pour chaque événement \\
Registre (lot Kroll) & 2,0\,s & 7,6 à 13,6\,s & toutes les ruches rejouées, tous les artefacts du registre \\
SRUM (21\,Mo) & 3,2\,s & 2,5\,s & 2 tables de plus, identités résolues \\
Prefetch (380) & 3,6\,s & 0,7\,s & heure locale datée, chemins avec lettre de volume \\
Raccourcis, jump lists & 1,8\,s & 0,3\,s & \\
Corbeille & 1,9\,s & 0,5\,s & \\
Amcache, ShimCache, shellbags & 2,0\,s & 0,5 à 1,4\,s & un seul passage pour les trois (clé registry) \\
\bottomrule
\end{longtable}

Sur les gros artefacts, \wac{} est plus rapide. Sur les petits, les outils EZ le
sont : toute conversion \wac{} paie un coût fixe de 1,5 à 3,4\,s, mesuré avec un
artefact trivial. Ce coût est celui de la copie de travail vérifiée des ruches, du
rejeu de leurs journaux, et de la lecture du fuseau et des profils du suspect.
C'est ce qui permet à \wac{} de dater chaque heure locale avec le décalage en
vigueur à sa date, et de donner la lettre de volume ; les outils EZ n'en
donnent pas.

\section{Collecte}

\begin{longtable}{@{}L{5.4cm}L{2.2cm}L{2.4cm}L{5.0cm}@{}}
\toprule
\textbf{Outil, réglage} & \textbf{Durée} & \textbf{Pièces} & \textbf{Volume} \\
\midrule\endhead
KAPE, \texttt{KapeTriage} & 58\,s & 1\,047 & 0,97\,Go \\
\wac{}, kit, sans binaires & 47\,s & 1\,898 & 2,54\,Go \\
\wac{}, kit, sans binaires, avant B-033 & 212\,s & 2\,622 & 1,69\,Go \\
\wac{}, kit complet (\texttt{binary: unverified}), avant B-033 & 454\,s & 61\,294 & 3,6\,Go \\
\bottomrule
\end{longtable}

Les deux premières lignes sont mesurées dans la même exécution, l'une après
l'autre. À périmètre comparable, \wac{} collecte désormais plus vite que KAPE,
en ramenant plus de deux fois plus de données. Il était cinq fois plus lent au
départ (fiche B-033) :
\begin{itemize}
  \item chaque lecture brute rouvrait son volume et reparcourait ses chemins
        depuis la racine ; les volumes sont maintenant ouverts une fois, avec le
        cache de leurs répertoires ;
  \item la collecte prenait les fichiers des 1\,426 fournisseurs d'événements
        déclarés (739\,Mo, 61\,s) ; elle ne prend plus que ceux des fournisseurs
        que citent les journaux collectés (324 sur la VM, 178\,Mo, 15\,s), ce qui
        suffit à reconstituer le texte de chaque événement.
\end{itemize}
Le kit complet vérifie en outre la signature Authenticode de chaque exécutable
de la machine et copie ceux dont la signature ne tient pas, ce que KAPE ne fait
pas.

% ============================================================ 4
\chapter{Précision}

\section{En faveur de \wac{}}

\begin{description}[style=nextline,leftmargin=1em]
  \item[Journaux de ruche rejoués dans l'ordre] Les journaux \fichier{.LOG1} et
        \fichier{.LOG2} gardent, derrière les entrées récentes, les restes d'une
        génération précédente. RECmd les applique dans l'ordre du fichier, donc
        après les plus récentes. Il ramène ainsi SYSTEM, SECURITY, DEFAULT et une
        ruche de profil de service à un état plus ancien : dans SYSTEM, la
        séquence 0x3618 au lieu de 0x3620, et quatre dernières exécutions BAM
        affichées 13 minutes trop tôt. \wac{} trie les entrées par séquence et
        s'arrête à la rupture de chaîne. Les journaux eux-mêmes le vérifient,
        sommes Marvin32 comprises.
  \item[Événements des journaux vidés] Un journal vidé garde ses enregistrements :
        \wac{} les relit et les marque \champ{EvtRecovered}, sans jamais les
        présenter comme vivants ; EvtxECmd ne les voit pas.
  \item[Valeurs d'événements comme Windows les rend] Hexadécimal \texttt{0x3e7},
        GUID entre accolades, booléens \texttt{true}, binaire \texttt{69060000},
        dates avec leur \texttt{Z}, marques U+200E conservées : les formes
        qu'affiche l'Observateur d'événements et que l'on recherche. EvtxECmd en
        change huit, et se trompe sur six valeurs, par exemple un booléen
        \texttt{False} que Windows rend \texttt{true}.
  \item[Rien d'inventé] Une valeur absente de la source n'est pas émise.
        AmcacheParser écrit 2\,340 valeurs par défaut absentes de la ruche
        (\texttt{IsPeFile: False} pour un exécutable PE), SrumECmd 7\,086
        (0 ou $-1$).
  \item[SHA-1 partiel de l'Amcache signalé] Pour un fichier de plus de 30\,Mio,
        Amcache ne hache que les 31\,457\,280 premiers octets (mesuré sur
        \fichier{mrt.exe}). \wac{} le dit (\champ{Sha1CoversBytes}) ; sans cela,
        la comparaison avec le SHA-1 du fichier entier ferait conclure à tort à
        une modification.
  \item[Pas de champ trompeur] L'indicateur « exécuté » du ShimCache, que donne
        AppCompatCacheParser, dit « non » pour 105 des 117 binaires dont le
        Prefetch prouve l'exécution. \wac{} ne le relève pas : ce « non » ferait
        conclure qu'un programme n'a jamais tourné.
  \item[Précision au 100\,ns] \wac{} donne chaque date à la précision de sa
        source. Par défaut, PECmd, LECmd et SrumECmd arrondissent à la seconde, et
        LECmd ignore \opt{mp} dans son CSV.
  \item[Heure locale datée] Chaque date est donnée en UTC et en heure locale du
        suspect, avec le décalage en vigueur à cette date. Les outils EZ ne
        donnent que l'UTC, sans le signaler par un \texttt{Z}.
  \item[Lectures d'autrui vérifiées] Les « FirstInteracted » et
        « LastInteracted » de SBECmd, le « CreatedOn » des comptes chez RECmd sont
        des lectures de dates d'écriture de clés. \wac{} publie ces dates
        elles-mêmes : chacun en tire la lecture qu'il veut, sans en hériter une
        toute faite. SBECmd donne en outre un \texttt{NodeSlot} faux (0) aux
        dossiers de premier niveau.
  \item[Erreurs de lecture ESE de SrumECmd] 17 lignes d'énergie sur 887 portent
        le contenu d'une autre ligne ou un identifiant absent de la table ;
        dissect.esedb donne raison à \wac{}.
  \item[Collecte] Chaque pièce de \wac{} porte MD5, SHA-1 et SHA-256, les huit
        dates NTFS de son original et son entrée MFT, dans un manifeste scellé.
        KAPE donne un SHA-1 et trois dates.
\end{description}

\section{En faveur des outils EZ}\label{sec:couverture}

\begin{description}[style=nextline,leftmargin=1em]
  \item[Couverture du registre] Le lot Kroll de RECmd lit des clés que \wac{}
        ignore : programmes installés, exclusions de Defender, RecentDocs,
        MountPoints2, réseaux connus, état de la mise à jour des dates d'accès
        NTFS, heure d'arrêt, entre autres. Fiche B-031.
  \item[Couverture de fichiers] KapeTriage ne copie plus, face à \wac{}, que
        les navigateurs et la chronologie Windows, prévus (B-016, B-010), et les
        \fichier{desktop.ini}, écartés. Les raccourcis hors de Recent, la ruche
        du profil Default, \fichier{\$Secure:\$SDS}, \fichier{\$Boot},
        \fichier{\$UsnJrnl:\$Max}, \fichier{\$TxfLog}, les traces WPP de
        Defender et les clés maîtresses DPAPI, qu'il était seul à prendre, sont
        collectés et convertis depuis la fiche B-032.
  \item[Vitesse sur les petits artefacts et en collecte] Voir le chapitre
        précédent.
  \item[Interprétations prêtes à lire] Les outils EZ découpent et nomment certains
        champs : exécutable, date et description d'une application SRUM, noms des
        dossiers connus, extension d'un fichier. \wac{} donne la valeur brute,
        telle que Windows la stocke.
\end{description}

\begin{attention}
Ce que \wac{} ne collecte pas, il ne peut pas le mettre en balance : la
couverture de KAPE reste plus large sur les navigateurs et plusieurs clés de
registre, jusqu'à ce que B-016, B-031 et B-032 soient faits.
\end{attention}

% ============================================================ 5
\chapter{Formats de sortie}

\section{Ce que produit chaque outil}

\begin{longtable}{@{}L{4.4cm}L{5.4cm}L{5.2cm}@{}}
\toprule
\textbf{Outil} & \textbf{Formats} & \textbf{Remarque} \\
\midrule\endhead
AmcacheParser & CSV & 8 fichiers, un par table \\
AppCompatCacheParser & CSV & \\
RBCmd & CSV & \\
SrumECmd & CSV & 7 fichiers \\
EvtxECmd & CSV, JSON, XML & \\
JLECmd & CSV, JSON, HTML & \\
LECmd & CSV, JSON, HTML, XML & \\
MFTECmd & CSV, JSON, bodyfile & \\
PECmd & CSV, JSON, HTML & \\
RECmd & CSV, JSON & plus un CSV par greffon \\
SBECmd & CSV, JSON & \\
KAPE & copie des fichiers, journal de copie CSV & \\
\wac{} & JSON pour chaque artefact & un fichier par artefact, plus le manifeste scellé \\
\bottomrule
\end{longtable}

\section{Ce qui les distingue}

Quatre outils sur onze ne produisent que du CSV. Chez les autres, le JSON est la
transcription de la ligne CSV : une ligne de texte par enregistrement (JSON
Lines), à plat.

\begin{description}[style=nextline,leftmargin=1em]
  \item[Types] Chez les outils EZ, les nombres sont des chaînes : un compteur
        d'exécutions s'écrit \texttt{"RunCount":"458"}, un numéro
        d'enregistrement \texttt{"EventRecordId":"21572"}. \wac{} écrit des
        nombres, des booléens, des chaînes.
  \item[Structures] Chez les outils EZ, une liste devient une chaîne jointe par
        des virgules (les fichiers chargés d'un Prefetch, ses répertoires), et les
        données d'un événement sont un JSON échappé \emph{dans} une chaîne
        (\texttt{Payload}), à décoder une seconde fois. \wac{} imbrique : les
        volumes d'un Prefetch et leurs répertoires, les éléments d'un raccourci,
        les paires \champ{Name}/\champ{Value} d'un événement.
  \item[Dates] Chez les outils EZ, l'UTC n'est pas signalé
        (\texttt{2026-09-30 23:04:46}) et la précision dépend d'une option. \wac{}
        donne l'ISO 8601 complet, en UTC (\texttt{…Z}) et en heure locale avec son
        décalage (\texttt{…+02:00}).
  \item[Absence de valeur] Les outils EZ écrivent une valeur par défaut, ou une
        cellule vide. \wac{} n'émet pas le champ ; un artefact non collecté est
        déclaré tel (\texttt{NotRequested}, \texttt{NotCollected}), pour qu'une
        liste vide ne passe pas pour une absence de trace.
  \item[Provenance] Le manifeste de \wac{} relie chaque valeur à la pièce dont
        elle vient, avec ses empreintes et ses dates d'origine.
  \item[Console] Les outils EZ détaillent leur travail sur la console ; certains
        y écrivent le décodage lui-même quand aucun fichier n'est demandé. \wac{}
        n'y écrit que l'avancement : tout ce qu'il trouve est dans les fichiers.
\end{description}

\begin{note}
Pour un analyste, la différence se voit au premier traitement : un JSON de
\wac{} se charge tel quel dans un outil d'analyse ou dans WAE, sans rien
redécouper ; une sortie EZ se traite comme un tableur, dont il faut reconvertir
les nombres, redécouper les listes et redécoder les données d'événement.
\end{note}

% ============================================================ 6
\chapter{Ce que la confrontation a appris à \wac{}}\label{chap:defauts}

Chaque défaut ci-dessous produisait un JSON valide et faux, ou incomplet. Tous
sont corrigés, avec un contrôle qui les aurait pris.

\begin{longtable}{@{}L{3.0cm}L{12.0cm}@{}}
\toprule
\textbf{Artefact} & \textbf{Défaut corrigé} \\
\midrule\endhead
\$MFT & enregistrements d'extension périmés fusionnés avec la base \\
EVTX & flottants tronqués à 6 chiffres ; binaire en « 69 06 00 00 » ; entiers hexadécimaux et GUID au mauvais format ; ProcessingErrorData ignoré (21 événements sans données) ; événements des journaux vidés non lus \\
Raccourcis & URI collée au bloc qui la suit ; lecture décalée de 2 octets sans liste d'ID ; blocs ExtraData non lus (machine, suivi) ; taille de la cible non lue ; DROID nuls émis \\
Jump lists & nombre d'ouvertures et rang MRU non lus ; format de DestList version 1 non pris en charge ; nom d'application altéré \\
Shellbags & dossier perdu avec sa descendance ; date de la clé du dossier, rang et NodeSlot absents ; dossier « Accueil » sans nom \\
Amcache & 2 tables sur 14 et 4 à 6 valeurs lues : ni SHA-1, ni taille, ni pilote, ni périphérique \\
ShimCache & tabulations des applications empaquetées remplacées par des espaces \\
SRUM & dates FILETIME laissées en nombres ; texte UTF-16 lu comme de l'ANSI \\
Registre & BAM limité aux comptes ayant un profil (un compte supprimé perdait ses exécutions) ; durée de premier plan de UserAssist non lue \\
Harnais & un contrôle qui plantait interrompait tous les suivants sans échec visible \\
\bottomrule
\end{longtable}

% ============================================================ 7
\chapter{Limites de la comparaison}

\begin{itemize}
  \item Une seule machine, virtuelle, sous Windows 11 : certains artefacts y
        sont peu nombreux (3 raccourcis de Recent, 2 clés USB). Les formats sont
        les mêmes sur une machine réelle, mais les cas rares n'y sont pas tous
        représentés.
  \item Les deux collectes se suivent au lieu d'être simultanées : les fichiers
        écrits entre elles (\$MFT, journaux, Prefetch) diffèrent légitimement.
  \item Les outils EZ ont été lancés avec leurs options de précision maximale ;
        par défaut, ils perdent davantage.
  \item La comparaison porte sur les artefacts communs. Ce que seul l'un des
        outils lit n'a pas de référence ici : son exactitude repose sur les
        contrôles du dossier de validation.
\end{itemize}

\end{document}
